\section{Discussion}
The transfer-function analysis places a limit on what a better input can achieve.
Under the single measured output and fixed initial conditions, excitation cannot
separate the original model's equivalent physical parameter combinations.
Estimating the three transfer coefficients or measuring an additional independent
quantity is a different scientific task from narrowing a four-parameter local
covariance. The existing data include a second sensor, but its usefulness for
breaking this equivalence requires a model that connects that sensor to the
states; the second trace alone does not establish that connection.

Initialization is part of the model specification. A delay that appears estimable
under cold delay states may become unobservable under equilibrated delay states
and constant input. Future laboratory records should include the heater history
before the fitting window and an independent ambient measurement. These records
would reduce dependence on assumptions made at the first sample.

The remaining discussion will assess the completed resampling, held-out
prediction, numerical sensitivity, and prospective-design results. Conditional
precision under a chosen residual model must be considered alongside error on
experiments excluded from calibration.
